% GENERATED FILE. Edit canonical lessons, then run npm run generate:books.
% canonical-source-hash: fnv1a64:e1d3b438468a512a
% canonical-lessons: SA-C235-pidadhati, SA-C235-aropayati, SA-C235-calayati, SA-C235-alingati, SA-C235-cumbati

\chapter{Covers, Plants, Drives, Embraces, Kisses}
\label{ch:sa-covers-plants-drives-embraces-kisses}
\hlchaptermodality{235}

\begin{chapteropening}
\textbf{By the end of this chapter:} \emph{I can say covers, plants, drives, embraces and kisses.}

Complete the last lesson of chapter 235: I can say covers, plants, drives, embraces and kisses.
\end{chapteropening}

\section[\texorpdfstring{\sk{पिदधाति} (pidadhāti)}{पिदधाति (pidadhāti)}]{\sk{पिदधाति} (pidadhāti) — covers, closes}

\label{lesson:SA-C235-pidadhati}



\begin{quote}
Before the new one: say the Sanskrit for finishes, then the Sanskrit for hands over.
\end{quote}

\subsection*{You'll want to know: \sk{पिदधाति}}
\textbf{\sk{पिदधाति}} — \emph{pidadhāti} — \textquotedblleft{}covers, closes\textquotedblright{}.

\begin{grammarlens}[title={What you've built}]
One more word for this chapter.
\end{grammarlens}

\subsection*{Your turn}
\begin{itemize}
\raggedright
  \item \emph{Say it:} \emph{pidadhāti}
  \item \emph{Say it:} \emph{pidadhāti}, once more
  \item \emph{Say it:} say \emph{pidadhāti}
  \item \emph{Recall:} say the Sanskrit for finishes, then the Sanskrit for hands over, then say \emph{pidadhāti} again
\end{itemize}

\begin{tcolorbox}[breakable,colback=teal!4,colframe=teal!35!black,title={Before you move on}]
What does \sk{पिदधाति} mean? (\textquotedblleft{}Covers, closes\textquotedblright{}.) Say it once more.
\end{tcolorbox}
\section[\texorpdfstring{\sk{आरोपयति} (āropayati)}{आरोपयति (āropayati)}]{\sk{आरोपयति} (āropayati) — plants, raises}

\label{lesson:SA-C235-aropayati}



\begin{quote}
Before the new one: say the Sanskrit for hands over, then the Sanskrit for covers.
\end{quote}

\subsection*{You'll want to know: \sk{आरोपयति}}
\textbf{\sk{आरोपयति}} — \emph{āropayati} — \textquotedblleft{}plants, raises\textquotedblright{}.

\begin{grammarlens}[title={What you've built}]
One more word for this chapter.
\end{grammarlens}

\subsection*{Your turn}
\begin{itemize}
\raggedright
  \item \emph{Say it:} \emph{āropayati}
  \item \emph{Say it:} \emph{āropayati}, once more
  \item \emph{Say it:} say \emph{āropayati}
  \item \emph{Recall:} say the Sanskrit for hands over, then the Sanskrit for covers, then say \emph{āropayati} again
\end{itemize}

\begin{tcolorbox}[breakable,colback=teal!4,colframe=teal!35!black,title={Before you move on}]
What does \sk{आरोपयति} mean? (\textquotedblleft{}Plants, raises\textquotedblright{}.) Say it once more.
\end{tcolorbox}
\section[\texorpdfstring{\sk{चालयति} (cālayati)}{चालयति (cālayati)}]{\sk{चालयति} (cālayati) — drives, moves}

\label{lesson:SA-C235-calayati}



\begin{quote}
Before the new one: say the Sanskrit for covers, then the Sanskrit for plants.
\end{quote}

\subsection*{You'll want to know: \sk{चालयति}}
\textbf{\sk{चालयति}} — \emph{cālayati} — \textquotedblleft{}drives, moves\textquotedblright{}.

\begin{grammarlens}[title={What you've built}]
One more word for this chapter.
\end{grammarlens}

\subsection*{Your turn}
\begin{itemize}
\raggedright
  \item \emph{Say it:} \emph{cālayati}
  \item \emph{Say it:} \emph{cālayati}, once more
  \item \emph{Say it:} say \emph{cālayati}
  \item \emph{Recall:} say the Sanskrit for covers, then the Sanskrit for plants, then say \emph{cālayati} again
\end{itemize}

\begin{tcolorbox}[breakable,colback=teal!4,colframe=teal!35!black,title={Before you move on}]
What does \sk{चालयति} mean? (\textquotedblleft{}Drives, moves\textquotedblright{}.) Say it once more.
\end{tcolorbox}
\section[\texorpdfstring{\sk{आलिंगति} (āliṅgati)}{आलिंगति (āliṅgati)}]{\sk{आलिंगति} (āliṅgati) — embraces}

\label{lesson:SA-C235-alingati}



\begin{quote}
Before the new one: say the Sanskrit for plants, then the Sanskrit for drives.
\end{quote}

\subsection*{You'll want to know: \sk{आलिंगति}}
\textbf{\sk{आलिंगति}} — \emph{āliṅgati} — \textquotedblleft{}embraces\textquotedblright{}.

\begin{grammarlens}[title={What you've built}]
One more word for this chapter.
\end{grammarlens}

\subsection*{Your turn}
\begin{itemize}
\raggedright
  \item \emph{Say it:} \emph{āliṅgati}
  \item \emph{Say it:} \emph{āliṅgati}, once more
  \item \emph{Say it:} say \emph{āliṅgati}
  \item \emph{Recall:} say the Sanskrit for plants, then the Sanskrit for drives, then say \emph{āliṅgati} again
\end{itemize}

\begin{tcolorbox}[breakable,colback=teal!4,colframe=teal!35!black,title={Before you move on}]
What does \sk{आलिंगति} mean? (\textquotedblleft{}Embraces\textquotedblright{}.) Say it once more.
\end{tcolorbox}
\section[\texorpdfstring{\sk{चुम्बति} (cumbati)}{चुम्बति (cumbati)}]{\sk{चुम्बति} (cumbati) — kisses}

\label{lesson:SA-C235-cumbati}



\begin{quote}
Before the new one: say the Sanskrit for drives, then the Sanskrit for embraces.
\end{quote}

\subsection*{You'll want to know: \sk{चुम्बति}}
\textbf{\sk{चुम्बति}} — \emph{cumbati} — \textquotedblleft{}kisses\textquotedblright{}.

\begin{grammarlens}[title={What you've built}]
One more word for this chapter.
\end{grammarlens}

\subsection*{Your turn}
\begin{itemize}
\raggedright
  \item \emph{Say it:} \emph{cumbati}
  \item \emph{Say it:} \emph{cumbati}, once more
  \item \emph{Say it:} say \emph{cumbati}
  \item \emph{Recall:} say the Sanskrit for drives, then the Sanskrit for embraces, then say \emph{cumbati} again
\end{itemize}

\begin{tcolorbox}[breakable,colback=teal!4,colframe=teal!35!black,title={Before you move on}]
What does \sk{चुम्बति} mean? (\textquotedblleft{}Kisses\textquotedblright{}.) Say it once more.
\end{tcolorbox}
